\documentclass[10pt,a4paper]{amsart}
\usepackage[margin=19mm]{geometry}
\usepackage{amsmath,amssymb,mathtools}
\usepackage[scr=rsfs,cal=euler]{mathalfa}
\usepackage{xfrac}
\usepackage[unicode,hidelinks,bookmarks=false]{hyperref}
\newcommand{\ie}{i.e.\ }
\newcommand{\cf}{cf.\ }
\newcommand{\resp}{respectively\ }
\newcommand{\Subsection}{\subsection}
\providecommand{\nobreakdash}{\nobreak\hbox{-}}
\setlength{\emergencystretch}{2em}
\multlinegap0pt
\usepackage{amssymb}
\usepackage{caption}
\usepackage{tikz}
\usepackage[shortlabels]{enumitem}
\usepackage{cancel}
\newtheorem{proposition}{Proposition}
\newcommand{\He}{\mathbb{H}}
\newcommand{\eT}{T_{\eps}}
\newcommand{\egrad}{\nabla_{\eps}}
\newcommand{\eps}{\varepsilon}
\newcommand{\escal}[2]{\langle {#1},{#2}\rangle_{\eps}}
\newcommand{\hLap}{\Delta_{\mathnormal{H}}}
\newcommand{\hgrad}{\nabla_{\mathnormal{H}}}
\newcommand{\hnorm}[1]{\|{#1}\|_{\mathnormal{H}}}
\newcommand{\scal}[2]{\langle {#1},{#2}\rangle}
\title{Global weak solutions for the inverse mean curvature flow in the Heisenberg group}
\author{Adriano Pisante, Eugenio Vecchi}
\date{Source: arXiv:2406.15123v3 (CC BY 4.0)}
\begin{document}
\maketitle

\noindent\emph{Source and licence.} Global weak solutions for the inverse mean curvature flow in the Heisenberg group. Version arXiv:2406.15123v3 (CC BY 4.0), distributed by arXiv under the Creative Commons Attribution 4.0 International licence, http://creativecommons.org/licenses/by/4.0/, as recorded in the arXiv licence metadata for this exact version and shown on the abstract page https://arxiv.org/abs/2406.15123v3. Full licence notice: \texttt{licenses/arxiv-2406.15123-CC-BY-4.0.txt}.\par\smallskip

\noindent\textit{Editorial scope.} Counted unit: the article's proposition SubBarrier (source0.tex lines 814--819). The article's complete original proof of it is delivered with it (source0.tex lines 820--961, one \texttt{proof} environment of 142 lines). No mathematics is written, repaired or re-derived: the statement and the proof are the article's own lines, in the article's own notation, and no retained line is substituted or re-flowed.  Retained uncounted as the source-local context the proof needs: lines 610--612; lines 657--664; lines 696--699; lines 727--730; lines 742--749; lines 751--754; lines 781--783.  Nothing was omitted from any retained span, so the package carries no omission record.  No retained line contains a citation, so the wrapper carries no bibliography and no bibliography entry.  No retained line contains a figure environment, an image-inclusion command or drawing code, so no image is delivered and the package declares no asset dependency.  Editorial formatting only, in addition: an AMS \texttt{amsart} preamble that restates the article's own declarations and macros used by the retained lines, namely the environments \texttt{proposition} on separate counters, together with the standard packages those lines need.  The retained environments print in this document as Proposition 1 in order of appearance, whereas the article prints the same items with the numbering of its own full document.  The complete native source of the article as distributed in the arXiv e-print for this exact version is bundled unchanged as \texttt{sources/161137/source0.tex}.
\begin{equation}\label{N}
N(x,y,t):= \| g\|^4= \left(|x|^2 + |y|^2\right)^2 + 16 t^2 = |z|^4 +16t^2 \, .
\end{equation}

\begin{align}
\label{XiPhi} X_i \Phi_{\alpha} &= \alpha N^{\alpha -1} X_i N = \alpha N^{\alpha-1} (4x_i |z|^2 - 16y_i t),\\
\label{YiPhi} Y_i \Phi_{\alpha} &= \alpha N^{\alpha -1} Y_i N = \alpha N^{\alpha-1} (4y_i |z|^2 + 16x_i t),\\
\label{TPhi} \eT \Phi_{\alpha} &= \alpha N^{\alpha -1} \eT N = \alpha N^{\alpha-1} 32 \eps t,\\
\eT X_i \Phi_{\alpha} &= \eps \left( \alpha (\alpha-1) N^{\alpha-2}(128 t (x_i |z|^2- 4y_i t)) - 16 \alpha N^{\alpha-1}y_i \right),\\
\eT Y_i \Phi_{\alpha} &= \eps \left( \alpha (\alpha-1) N^{\alpha-2}(128 t (y_i |z|^2+ 4x_i t)) + 16 \alpha N^{\alpha-1}x_i \right),\\
\label{TTPhi} \eT \eT \Phi_{\alpha} &= \eps^2 \left( \alpha (\alpha-1)N^{\alpha-2} (32t)^2 + 32 \alpha N^{\alpha-1}\right).
\end{align}

\begin{equation}
\label{nabla0Phi2}
\nabla_0 \Phi = \alpha N^{\alpha -1} \nabla_0 N \quad \textrm{and } \quad |\nabla_0 \Phi|_0^2 = 16 \alpha^2 |z|^2 N^{2\alpha -1} , 
\end{equation}

\begin{equation}
\label{firstpdec}
\frac{p-2}2\scal{\nabla_0\left( |\nabla_0 \Phi|_0^{2}\right)}{\nabla_0 \Phi}_0 = (4 \alpha)^3 (4\alpha-1) (p-2) |z|^4  N^{3\alpha -2} \, .
\end{equation}

\begin{equation}
\label{Delta0Phi}
\begin{aligned}
\Delta_0 \Phi &= \alpha (\alpha-1)N^{\alpha-2}|\nabla_0 N|_0^2 + \alpha N^{\alpha-1} \Delta_0 N \\
&= 16 \alpha (\alpha-1)|z|^2 N^{\alpha-1} + 8(2+n) \alpha |z|^2 N^{\alpha -1} \\
&= 8\alpha (2\alpha +n) |z|^2 \, N^{\alpha-1} \, , 
\end{aligned}
\end{equation}

\begin{equation}
\label{secondpdec}
|\nabla_0 \Phi|_0^{2}\Delta_0 \Phi = (4 \alpha)^{3}2  (2\alpha +n) |z|^4 N^{3\alpha-2}.
\end{equation}

\begin{equation}\label{PhiR0}
\Phi^{(g_0, R_0)}_{\alpha}(x,y,t):= \dfrac{\Phi_{\alpha} \circ L_{g_0^{-1}}(x,y,t)}{R_{0}^{4\alpha}} \, .
\end{equation}

\begin{proposition}\label{SubBarrier}
For any $R_0>0$, $1<p<Q$ and $\varepsilon \in (0,1]$ let $K\geq \frac{Q-p}4 + \varepsilon^4 \left(\frac{2}{R_0^2} +  \frac{8 Q}{R_0^4} \right) $. Then for any $g_0 \in \mathbb{H}^n$ the function $\Phi^{(g_0, R_0)}_{\alpha}$ defined in \eqref{PhiR0} with $\alpha (p-1) = -K$ satisfies
\begin{equation}\label{goal}
\mathrm{div}_{\eps}\left(  |\nabla_{\eps} \Phi^{(g_0, R_0)}_{\alpha}|_{\eps}^{(p-2)}\nabla_{\eps}\Phi^{(g_0, R_0)}_{\alpha} \right) \geq 0 \quad \textrm{in } B^{c}_{R_0}(g_0) \, .
\end{equation}
\end{proposition}

\begin{proof}
Without loss of generality, we can assume up to a translation that $g_0 = 0$. To simplify further the
readability, we set $\Phi=\Phi_{\alpha}$ and we define 
\begin{equation}\label{def:Phi'}
  \Phi':= \Phi^{(0, R_0)}_{\alpha} = R_{0}^{-4\alpha} \Phi \, .
  \end{equation}
We further observe that
\begin{equation}
\label{decpepslap}
\begin{aligned}
\mathrm{div}_{\eps}\left(  |\nabla_{\eps} \Phi'|_{\eps}^{(p-2)}\nabla_{\eps}\Phi' \right) &= R_{0}^{-4\alpha (p-1)} \mathrm{div}_{\eps}\left( \left( |\nabla_{\eps} \Phi|_{\eps}^{2}\right)^{(p-2)/2}\nabla_{\eps}\Phi \right) \\
&= R_{0}^{-4\alpha (p-1)} \left[\left(|\egrad \Phi|_\varepsilon^2 \right)^{\tfrac{p-2}{2}} \Delta_{\eps}\Phi + \escal{\egrad \left(|\egrad \Phi|_\varepsilon^{2} \right)^{\tfrac{p-2}{2}}}{\egrad \Phi}\right]\\
&=: R_{0}^{-4\alpha (p-1)} \left( (i) + (ii) \right).
\end{aligned}
\end{equation}
Now, for every smooth vector field $W$, it holds that
\begin{equation}
\label{Waction}
\begin{aligned} 
W &\left(|\egrad \Phi|_\varepsilon^2 \right)^{\tfrac{p-2}{2}} = \tfrac{p-2}{2} \left(|\egrad \Phi|_\varepsilon^2 \right)^{\tfrac{p-4}{2}} W (|\egrad \Phi|_\varepsilon^{2})\\
&= \tfrac{p-2}{2} \left(|\egrad \Phi|_\varepsilon^2 \right)^{\tfrac{p-4}{2}} W \left( \sum_{i=1}^{n}\left[(X_i \Phi)^2 + (Y_i \Phi)^2\right] + (\eT \Phi)^2 \right) \\
&= (p-2) \left(|\egrad \Phi|_\varepsilon^2 \right)^{\tfrac{p-4}{2}} \left[ \sum_{i=1}^{n} \left( X_i \Phi_\alpha WX_i \Phi + Y_i \Phi_\alpha WY_i \Phi \right) + \eT \Phi_\alpha W\eT \Phi \right].
\end{aligned}
\end{equation}
Applying \eqref{Waction} with $W=\nabla_\varepsilon \Phi$ in $B_{R_0}(0)^c$ yields 
\begin{equation*}
\begin{aligned}
(ii) &= (p-2) \left(|\egrad \Phi|_\varepsilon^2 \right)^{\tfrac{p-4}{2}} \sum_{j=1}^{n}\left[(X_j \Phi)\left(\sum_{i=1}^{n} (X_i\Phi X_j X_i \Phi + Y_i \Phi X_j Y_i \Phi) + \eT \Phi X_j \eT \Phi \right)\right]\\
&+ (p-2) \left(|\egrad \Phi|_\varepsilon^2 \right)^{\tfrac{p-4}{2}}\sum_{j=1}^{n}\left[(Y_j \Phi) \left(\sum_{i=1}^{n}(X_i \Phi Y_j X_i \Phi + Y_i \Phi Y_j Y_i \Phi )+ \eT \Phi Y_j \eT \Phi\right)\right]\\
&+ (p-2) \left(|\egrad \Phi|_\varepsilon^2 \right)^{\tfrac{p-4}{2}} \eT \Phi \left( \sum_{i=1}^{n}(X_i \Phi \eT X_i \Phi + Y_i \Phi \eT Y_i \Phi) + \eT \Phi \eT \eT \Phi \right).
\end{aligned}
\end{equation*}
Commuting $X_i$ and $Y_i$ with $T_\varepsilon=\varepsilon T$ and reordering with respect to the powers of $\eps$ we obtain
\begin{equation}
\label{(ii)}
\begin{aligned}
(ii)&= (p-2)\left(|\egrad \Phi|_\varepsilon^2 \right)^{\tfrac{p-4}{2}} \sum_{j=1}^{n}\left[(X_j \Phi)\sum_{i=1}^{n} (X_i\Phi X_j X_i \Phi + Y_i \Phi X_j Y_i \Phi) \right]\\
&+ (p-2) \left(|\egrad \Phi|_\varepsilon^2 \right)^{\tfrac{p-4}{2}}\sum_{j=1}^{n}\left[ (Y_j \Phi) \sum_{i=1}^{n}(X_i \Phi Y_j X_i \Phi + Y_i \Phi Y_j Y_i \Phi ) \right]\\
&+2\eps^2 (p-2)\left(|\egrad \Phi|_\varepsilon^2 \right)^{\tfrac{p-4}{2}} \sum_{j=1}^{n}\left(X_j \Phi T\Phi  X_j T\Phi + Y_j \Phi T\Phi Y_j T\Phi \right) \\
&+ \eps^4 (p-2) \left(|\egrad \Phi|_\varepsilon^2 \right)^{\tfrac{p-4}{2}}  (T\Phi )^2 TT\Phi \\
&=  (p-2)\left(|\egrad \Phi|_\varepsilon^2 \right)^{\tfrac{p-4}{2}} \left[  \frac12\scal{\nabla_0 |\nabla_0 \Phi|^2_0}{\nabla_0 \Phi}_0 +\varepsilon^2 T\Phi T \left( |\nabla_0 \Phi|^2_0\right)+ \varepsilon^4(T\Phi)^2 TT\Phi \right] \,.
\end{aligned}
\end{equation}
On the other hand, using again the identity $T_\varepsilon= \varepsilon T$ we also have 
\begin{equation}
\label{(i)}
\begin{aligned}
(i) &= \left(|\egrad \Phi|_\varepsilon^2 \right)^{\tfrac{p-4}{2}} \left( |\nabla_0 \Phi|_0^{2} + \eps^2 (T\Phi)^2 \right) (\Delta_0 \Phi + \eps^2 TT\Phi)\\
&= \left(|\egrad \Phi|_\varepsilon^2 \right)^{\tfrac{p-4}{2}} \left[ |\nabla_0 \Phi|_0^{2}\Delta_0 \Phi+ \varepsilon^2 \left( |\nabla_0 \Phi|_0^{2}TT\Phi+ (T\Phi)^2 \Delta_0 \Phi\right) +\varepsilon^4 (T\Phi)^2TT\Phi \right]\,.
\end{aligned}
\end{equation}

Thus, \eqref{decpepslap}, \eqref{(i)} and \eqref{(ii)} yield 
\begin{equation}
\label{decpepslap2}
\mathrm{div}_{\eps} \left( |\nabla_{\eps} \Phi|_{\eps}^{(p-2)}\nabla_{\eps}\Phi \right) =(i)+(ii)= \left(|\nabla_{\eps}\Phi|_{\eps}^2 \right)^{\tfrac{p-4}{2}}\left[ (I) + (II) + (III)  \right],
\end{equation}
\noindent where
\begin{equation}
\label{1-2-3}
\begin{aligned}
(I):&= |\nabla_0 \Phi|_0^{2}\Delta_0 \Phi+\frac{(p-2)}2   \scal{\nabla_0 |\nabla_0 \Phi|^2_0}{\nabla_0 \Phi}_0 \, ,\\
(II):&= \varepsilon^2 \left( |\nabla_\varepsilon \Phi|_0^{2}TT\Phi+ (T\Phi)^2 \Delta_0 \Phi    +(p-2)T\Phi T \left( |\nabla_0 \Phi|^2_0\right) \right) \, , \\
(III):&= \eps^4 (p-1)  (T\Phi)^2 (TT\Phi) \, .
\end{aligned}
\end{equation}

%\begin{align*}
%(I):&= (\hnorm{\hgrad v}^{2} + \eps (Tv)^2) (\hLap v + \eps TTv)\\
 %  &= \hnorm{\hgrad v}^{2} \hLap v +  \eps (Tv)^2 \hLap v + \eps \hnorm{\hgrad v}^2 TTv + \eps^2 (Tv)^2 TTv,\\
  % (II):&= (p-2) \sum_{j=1}^{n} \left[ (X_j v) \sum_{i=1}^{n}\left( X_i v X_j X_i v + Y_i v X_j Y_i v\right) + Y_j v \sum_{i=1}^{n} \left(X_i v Y_j X_i v + Y_i v Y_j Y_i v \right)\right],\\
%(III):&= 2\eps(p-2)  \sum_{j=1}^{n}\left(X_j v Tv X_j Tv + Y_j v Tv Y_j Tv \right),\\
%(IV):&= \eps^2 (p-2)  (Tv)^2 (TTv).
%\end{align*}
\noindent Concerning the first term in \eqref{1-2-3} identities \eqref{firstpdec} and \eqref{secondpdec} give
\begin{equation}
\label{ineq(1)}
(I)= (4\alpha)^3 |z|^4 N^{3\alpha-2} (4\alpha (p-1)+Q-p)
%\geq  (4 \alpha)^3|z|^4 N^{3\alpha-2} (-4K+Q-1) 
>0 \, ,
\end{equation}
\noindent for any $K>\frac{Q-p}4$.

Next for $\alpha<0$ and restricting to $\He^n\setminus B_{R_0}$ we bound separately from below the terms $(II)$ and $(III)$ in \eqref{1-2-3}.
To this end, notice that on $\He^n\setminus B_{R_0}(0)$ we have $N \geq R_0^4$ an therefore 
\begin{equation}\label{StimaNR0}
-N^{\gamma} \geq -R_0^{4\gamma}, \quad \textrm{for every } \gamma<0.
\end{equation}

%\begin{align*}
%(II) &= \eps \left[ (Tv)^2 \hLap v +  \hnorm{\hgrad v}^2 \, TTv + 2(p-2)(Tv) \sum_{j=1}^{n}(X_jv X_j Tv + Y_j v Y_j Tv)\right].\\
%(III) &= \eps^2 (p-1) (Tv)^2 (TTv) = \eps^2 (p-1) (32 \alpha N^{\alpha-1}t)^2 \left[32 \alpha N^{\alpha-2} (N+(\alpha-1)32t^2)\right]\\
%&= \eps^2 (32 \alpha)^3 (p-1) N^{3\alpha -4} t^2 \left[ |z|^4 + \underbrace{16 (2\alpha -1)t^2}_{\leq 0} \right], 
%\end{align*}
%We stress that $(A)$ contains all the terms in $(I)$, $(II)$,$(III)$ and $(IV)$
%with $\eps^2$, while $(B)$ all those terms with $\eps$.
\noindent Concerning $(II)$, using \eqref{TPhi}, \eqref{TTPhi}, \eqref{nabla0Phi2} and \eqref{Delta0Phi} we compute all the terms and obtain
\begin{align*}
(T\Phi)^2 \Delta_0 \Phi &= 16 (8\alpha)^3 (2\alpha +n) N^{3\alpha -3}\, |z|^2 \, t^2,\\
|\nabla_0 \Phi|^2 TT\Phi &= (8 \alpha)^3  N^{3\alpha -3}\, |z|^2 (|z|^4 + (2\alpha-1)16t^2) \, ,\\
(p-2)T\Phi T \left( |\nabla_0 \Phi|^2_0\right)&= 32  (p-2) (8\alpha)^3 (2\alpha -1)N^{3\alpha -3} \, |z|^2 \, t^2 \, ,
\end{align*}  
whence
\begin{equation}
\label{eq:II}
\begin{aligned}(II)&=\varepsilon^2 (8 \alpha)^3 N^{3\alpha-3}|z|^2 \left[ 16 (2\alpha+n) t^2+ |z|^4+16(2\alpha-1)t^2 +32(p-2)(2\alpha-1) t^2\right]\\
&=\varepsilon^2 (8 \alpha)^3 N^{3\alpha-3}|z|^2 \left[ |z|^4+ 8t^2(Q+4(2\alpha-1)(p-1))\right] \, .
\end{aligned}
\end{equation}
Note that for $p>1$ we have $Q+4(2\alpha-1)(p-1)<0$ by our choice of $\alpha$, hence combining on $\He^n\setminus B_{R_0}$ the last identity with \eqref{N} and \eqref{StimaNR0} we obtain uniformly w.r.to $\varepsilon \in (0,1]$
\begin{equation}
\label{ineq(2)}
(II) \geq  \varepsilon^2 (8 \alpha)^3 |z|^6 N^{3\alpha-3}  \geq \varepsilon^4 (4 \alpha)^3 |z|^4 N^{3\alpha-2} 8 N^{-1/2} \geq \varepsilon^4 (4 \alpha)^3 |z|^4 N^{3\alpha-2} 8 R_0^{-2}\, .
\end{equation} 

%Summing all together the terms with $|z|^2 t^2$, we find
%\begin{equation*}
%\begin{aligned}
%2 &\cdot 16^3 \alpha^3 N^{3\alpha-3}|z|^2 t^2 \left( 2\alpha + n +2(p-2)(2\alpha -1) + 2\alpha -1 \right)\\
%&= 2 \cdot 16^3 \alpha^3 N^{3\alpha-3}|z|^2 t^2 (4\alpha + n-1 + 2(p-2)(2\alpha -1)) \geq 0.
%\end{aligned}
%\end{equation*}
%Concerning $(B)$, we are left with only with
%the term $512 \eps \alpha^3 N^{3\alpha-3}|z|^6$. \\
\noindent On the other hand, applying \eqref{TPhi} and \eqref{TTPhi} we get
\begin{equation}
\label{eq:III}
\begin{aligned}(III) &= \eps^4 (p-1) (T\Phi)^2 (TT\Phi) = \eps^4 (p-1) (32 \alpha N^{\alpha-1}t)^2 \left[32 \alpha N^{\alpha-2} (N+(\alpha-1)32t^2)\right]\\
&= \eps^4 (32 \alpha)^3 (p-1) N^{3\alpha -4} t^2 \left[ |z|^4 + \underbrace{16 (2\alpha -1)t^2}_{\leq 0} \right], 
\end{aligned}
\end{equation}
therefore on $\He^n\setminus B_{R_0}$ for all $\varepsilon \in (0,1]$ by \eqref{N} and \eqref{StimaNR0} we obtain
\begin{equation}
\label{ineq(3)}
(III)\geq \varepsilon^4 8^3 (4\alpha)^3 (p-1) N^{3\alpha-4} t^2 |z|^4\geq \varepsilon^4 32 (4\alpha)^3 (p-1)  |z|^4 N^{3\alpha-2}   N^{-1}\geq  \varepsilon^4 (4\alpha)^3  |z|^4 N^{3\alpha-2} 32 Q   R_0^{-4}\,.
\end{equation}

Combining \eqref{decpepslap2}, \eqref{ineq(1)}, \eqref{ineq(2)}, and \eqref{ineq(3)} in the whole $\He^n\setminus B_{R_0}$ we infer

 \[ \mathrm{div}_{\eps} \left( |\nabla_{\eps} \Phi|_{\eps}^{(p-2)}\nabla_{\eps}\Phi \right) \geq(4\alpha^3) |z|^4 N^{3\alpha-2} \left(|\nabla_{\eps}\Phi|_{\eps}^2 \right)^{\tfrac{p-4}{2}}\left[ -4K+Q-p + \varepsilon^4 \left(\frac{8}{R_0^2} +  \frac{32 Q}{R_0^4} \right) \right] \, ,\]
 whence the conclusion follows from the choice of $K$ as $\alpha<0$.
 \end{proof}

\end{document}
